El primer tipo de ecuaci\'on de segundo orden que revisaremos ser\'a la ecuaci\'on homog\'enea de coeficientes constantes, la cual es:
\[
ay''+by'+cy=0,
\]
donde $a,b,c\in\mathbb{R}$ y $a\neq 0$. Este tipo de ecuaciones emerge de manera natural en problemas del tipo masa-resorte. La ecuaci\'on se dice homog\'enea, ya que no hay una funci\'on externa (fuerza) $f(t)$. En el caso en que existe una fuerza externa, la ecuaci\'on de coeficientes constantes se convierte en:
\[
ay''+by'+cy=f(t).
\]

La idea para resolver la ecuaci\'on homog\'enea se la debemos nuevamente al gran L.~Euler y se trata de lo siguiente: buscamos una funci\'on $y$ para la cual su segunda derivada sea una combinaci\'on de la funci\'on buscada y de su primera derivada. Esta simple idea sugiere pensar en funciones del tipo $y=e^{rt}$, ya que sus derivadas son m\'ultiplos de ella misma. De hecho, si reemplazamos:
\[
ar^2e^{rt}+bre^{rt}+ce^{rt}=0\Longrightarrow e^{rt}(ar^2+br+c)=0.
\]
Como la funci\'on $e^{rt}$ nunca es cero, tenemos que
\[
ar^2+br+c=0.
\]
A esta ecuaci\'on la llamaremos la \textbf{ecuaci\'on auxiliar} (o caracter\'istica) de la ecuaci\'on homog\'enea de coeficientes constantes.

\begin{example}
  Consideremos la ecuaci\'on dada por:
  \[
  y''+5y'-6y=0.
  \]
  Su respectiva ecuaci\'on auxiliar es:
  \[
  r^2+5r-6=(r-1)(r+6)=0,
  \]
  por lo tanto, las funciones $e^{t}$ y $e^{-6t}$ son soluciones de la ecuaci\'on dada.
\end{example}
\newpage
%[Principio de superposici\'on]
\begin{remark}
    Notemos que la funci\'on $y(t)\equiv 0$ siempre es soluci\'on de la ecuaci\'on homog\'enea. M\'as a\'un, dadas dos soluciones $y_1,y_2$ de ella y dos constantes $c_1,c_2\in \mathbb{R}$, la funci\'on
    \[
    y(t)=c_1y_1(t)+c_2y_2(t)
    \]
    es tambi\'en soluci\'on. Esto es f\'acil de verificar reemplazando en la ecuaci\'on homog\'enea y usando la linealidad de la derivada:
    \begin{align*}
      &a(c_1y_1+c_2y_2)''+b(c_1y_1+c_2y_2)'+c(c_1y_1+c_2y_2)\\
      &\qquad= a(c_1y_1''+c_2y_2'')+b(c_1y_1'+c_2y_2')+c(c_1y_1+c_2y_2)\\
      &\qquad= c_1\underbrace{(ay_1''+by_1'+cy_1)}_{=0}+c_2\underbrace{(ay_2''+by_2'+cy_2)}_{=0}=0.
    \end{align*}
    Los dos grados de libertad que entregan las constantes $c_1,c_2$ nos sugieren que, a diferencia de las ecuaciones de orden uno, las ecuaciones de segundo orden requieren m\'as condiciones iniciales para poder determinar soluciones particulares.

    En t\'erminos de un sistema mec\'anico, no nos basta con la posici\'on inicial $y(0)$: tambi\'en requerimos la velocidad inicial del sistema $y'(0)$ para determinar su comportamiento. A modo de ejemplo, un PVI de segundo orden se ve como a continuaci\'on.
\end{remark}

\begin{example}
     Considere el PVI dado por
     \[
     y''+2y'-y=0,\qquad y(0)=0,\qquad y'(0)=-1.
     \]
     Al igual que en el ejemplo anterior, utilizamos la ecuaci\'on auxiliar para encontrar las soluciones de la forma $y(t)=e^{rt}$. En este caso,
     \[
     r^2+2r-1=0\Longrightarrow r_1=-1+\sqrt{2},\qquad r_2=-1-\sqrt{2}.
     \]
     Con lo cual la soluci\'on general ser\'a:
     \[
     y(t)=c_1e^{(-1+\sqrt{2})t}+c_2e^{(-1-\sqrt{2})t}.
     \]
     A continuaci\'on debemos determinar los valores de $c_1,c_2$ a partir de los valores iniciales dados, es decir:
     \begin{align*}
         y(0)&=c_1e^0+c_2e^0,\\
         y'(0)&=(-1+\sqrt{2})c_1e^0+(-1-\sqrt{2})c_2e^0,
     \end{align*}
     de lo anterior tenemos el sistema de ecuaciones:
     \begin{align*}
         0&=c_1+c_2,\\
         -1&=(-1+\sqrt{2})c_1+(-1-\sqrt{2})c_2.
     \end{align*}
     De la primera ecuaci\'on, $c_2=-c_1$; reemplazando en la segunda obtenemos $2\sqrt{2}\,c_1=-1$. Por lo tanto $c_1=-\dfrac{\sqrt{2}}{4}$, $c_2=\dfrac{\sqrt{2}}{4}$, y la respectiva soluci\'on particular es:
     \[
     y(t)=-\dfrac{\sqrt{2}}{4}e^{(-1+\sqrt{2})t}+\dfrac{\sqrt{2}}{4}e^{(-1-\sqrt{2})t}.
     \]
\end{example}

En el cap\'itulo anterior revisamos un teorema que nos aseguraba, bajo ciertas condiciones, la existencia y unicidad de soluci\'on para las ecuaciones de orden uno. A continuaci\'on revisaremos el resultado equivalente para PVI con ecuaciones de orden dos.

\begin{theorem}[Existencia y unicidad]
    Para cualquier colecci\'on de n\'umeros reales $a,b,c,t_0,Y_0,Y_1$ con $a\neq 0$, existe una \'unica soluci\'on del problema de valor inicial
    \[
    ay''+by'+cy=0,\qquad y(t_0)=Y_0,\qquad y'(t_0)=Y_1,
    \]
    y esta soluci\'on es v\'alida para todo $t\in\, ]-\infty,\infty[$.
\end{theorem}

Como veremos a continuaci\'on, el hecho de que el teorema anterior nos asegure la unicidad de la soluci\'on es de gran ayuda: nos bastar\'a encontrar (de cualquier manera) dos soluciones $y_1,y_2$ que sean \emph{esencialmente distintas} para que la funci\'on $c_1y_1+c_2y_2$, con constantes adecuadas, sea la soluci\'on buscada. Notemos que no basta con que sean distintas: por ejemplo, $e^{t}$ y $2e^{t}$ son distintas, pero cualquier combinaci\'on de ellas es de la forma $Ce^{t}$, y con una sola constante no podemos satisfacer dos condiciones iniciales. Para precisar esta idea, definimos el siguiente concepto.

\begin{definition}
    Dos funciones $y_1(t),\,y_2(t)$ se dicen \textbf{linealmente independientes} (L.I.) sobre el intervalo $I\subseteq \mathbb{R}$ si la \'unica forma de tener
    \[
    c_1y_1(t)+c_2y_2(t)=0\quad\text{para todo } t\in I
    \]
    es con $c_1=c_2=0$. En caso contrario diremos que las funciones son \textbf{linealmente dependientes} (L.D.) sobre $I$.

    Para dos funciones, esto equivale a decir que ninguna de ellas es un m\'ultiplo constante de la otra en $I$; es decir, no existe una constante $k$ tal que $y_1(t)=k\,y_2(t)$ para todo $t\in I$ (ni $y_2(t)=k\,y_1(t)$ para todo $t\in I$).
\end{definition}

\begin{example}
    Las funciones
    \[
    y_1(t)=\sin(t)\cos(t),\qquad y_2(t)=\sin(2t)
    \]
    son linealmente dependientes en $\mathbb{R}$, ya que por la identidad del \'angulo doble $y_1(t)=\tfrac{1}{2}\,y_2(t)$ para todo $t$.

    En cambio, las funciones
    \[
    y_1(t)=e^{t},\qquad y_2(t)=e^{4t}
    \]
    son linealmente independientes en $\mathbb{R}$, pues el cociente $y_2(t)/y_1(t)=e^{3t}$ no es constante.
\end{example}

Para verificar la independencia lineal de manera sistem\'atica, introducimos la siguiente herramienta.

\begin{definition}
    Dadas dos funciones derivables $y_1,y_2$, se define el \textbf{Wronskiano} de ellas como:
    \[
    W[y_1,y_2](t)=\begin{vmatrix}
        y_1(t) & y_2(t)\\
        y_1'(t) & y_2'(t)
    \end{vmatrix}=y_1(t)y_2'(t)-y_1'(t)y_2(t).
    \]
\end{definition}

\begin{remark}
    Si existe al menos un punto $t_0\in I$ tal que $W[y_1,y_2](t_0)\neq 0$, entonces las funciones $y_1,y_2$ son L.I. en $I$.

    El rec\'iproco es falso para funciones arbitrarias: por ejemplo, $y_1(t)=t^2$ e $y_2(t)=t|t|$ tienen Wronskiano id\'enticamente nulo en $\mathbb{R}$ y, sin embargo, son L.I. en $\mathbb{R}$. No obstante, cuando $y_1,y_2$ son \emph{soluciones} de la ecuaci\'on homog\'enea, su Wronskiano o bien es id\'enticamente cero, o bien no se anula nunca (esto lo justificaremos m\'as adelante con la f\'ormula de Abel). En ese caso, las soluciones son L.I. si y solo si su Wronskiano es distinto de cero.
\end{remark}

\begin{example}
    Para las funciones $y_1(t)=e^{t}$, $y_2(t)=e^{4t}$ se tiene:
    \[
    W[y_1,y_2](t)=\begin{vmatrix}
        e^t & e^{4t}\\
        e^t & 4e^{4t}
    \end{vmatrix}=4e^{5t}-e^{5t}=3e^{5t}\neq 0,\qquad \forall\, t\in\, ]-\infty,\infty[,
    \]
    lo que confirma que son L.I.
\end{example}

El siguiente resultado relaciona la existencia de soluciones con los conceptos reci\'en introducidos.

\begin{theorem}
    Si $y_1(t),y_2(t)$ son dos soluciones L.I. en $]-\infty,\infty[$ de la ecuaci\'on de coeficientes constantes, entonces existen \'unicas constantes $c_1,c_2$ de manera que $c_1y_1(t)+c_2y_2(t)$ satisface el problema de valor inicial
    \[
    ay''+by'+cy=0,\qquad y(t_0)=Y_0,\qquad y'(t_0)=Y_1,
    \]
    en dicho intervalo. Por esta raz\'on, a la expresi\'on $y_h=c_1y_1+c_2y_2$ la llamamos \textbf{soluci\'on general} de la ecuaci\'on homog\'enea.
\end{theorem}

\begin{remark}
    Es instructivo ver d\'onde aparece el Wronskiano en el teorema anterior. Imponer las condiciones iniciales a $c_1y_1+c_2y_2$ lleva al sistema
    \begin{align*}
        c_1y_1(t_0)+c_2y_2(t_0)&=Y_0,\\
        c_1y_1'(t_0)+c_2y_2'(t_0)&=Y_1,
    \end{align*}
    cuyo determinante es exactamente $W[y_1,y_2](t_0)$. Por lo tanto, el sistema tiene soluci\'on \'unica $(c_1,c_2)$ precisamente cuando $W[y_1,y_2](t_0)\neq 0$, es decir, cuando las soluciones son L.I.
\end{remark}

Para la ecuaci\'on de coeficientes constantes tenemos, por lo tanto, tres casos seg\'un el tipo de ra\'ices de la ecuaci\'on auxiliar.

\subsection*{Caso 1: ra\'ices reales y distintas}

\begin{theorem}
    Si la ecuaci\'on auxiliar $ar^2+br+c=0$ posee dos soluciones reales y distintas $r_1,r_2$, entonces la ecuaci\'on homog\'enea tiene por soluci\'on general:
    \[
    y_h=c_1e^{r_1t}+c_2e^{r_2t}.
    \]
\end{theorem}

En efecto, $W[e^{r_1t},e^{r_2t}](t)=(r_2-r_1)e^{(r_1+r_2)t}\neq 0$, pues $r_1\neq r_2$. Este es el caso que encontramos en los dos primeros ejemplos de la secci\'on.

\subsection*{Caso 2: una ra\'iz real doble}

\begin{theorem}
    Si la ecuaci\'on auxiliar $ar^2+br+c=0$ posee una \'unica ra\'iz real $r$ (ra\'iz doble), entonces la ecuaci\'on homog\'enea tiene por soluci\'on general:
    \[
    y_h=c_1e^{rt}+c_2te^{rt}.
    \]
\end{theorem}

\begin{proof}[Justificaci\'on]
    Si $r$ es ra\'iz doble, entonces $r=-\dfrac{b}{2a}$, es decir, $2ar+b=0$. Para $y=te^{rt}$ tenemos $y'=(1+rt)e^{rt}$ e $y''=(2r+r^2t)e^{rt}$, luego
    \[
    ay''+by'+cy=e^{rt}\big[t\underbrace{(ar^2+br+c)}_{=0}+\underbrace{(2ar+b)}_{=0}\big]=0.
    \]
    Adem\'as, $W[e^{rt},te^{rt}](t)=e^{2rt}\neq 0$, por lo que ambas soluciones son L.I.
\end{proof}

\begin{example}
    Consideremos la ecuaci\'on
    \[
    y''+4y'+4y=0.
    \]
    Su ecuaci\'on auxiliar es $r^2+4r+4=(r+2)^2=0$, que tiene una \'unica ra\'iz $r=-2$. Con esto, por el resultado anterior,
    \[
    y_h(t)=c_1e^{-2t}+c_2te^{-2t}.
    \]
    Verifiquemos directamente que $y_2(t)=te^{-2t}$ es soluci\'on de nuestra ecuaci\'on:
    \begin{align*}
        y_2(t)&=te^{-2t},\\
        y_2'(t)&=e^{-2t}-2te^{-2t},\\
        y_2''(t)&=-4e^{-2t}+4te^{-2t},
    \end{align*}
    y reemplazando:
    \[
    y_2''+4y_2'+4y_2=\left(-4e^{-2t}+4te^{-2t}\right)+4\left(e^{-2t}-2te^{-2t}\right)+4te^{-2t}=0.
    \]
\end{example}

\subsection*{Caso 3: ra\'ices complejas conjugadas}

\begin{theorem}
    Si la ecuaci\'on auxiliar posee ra\'ices complejas $\alpha\pm i\beta$ (con $\beta\neq 0$), entonces dos soluciones linealmente independientes de la ecuaci\'on de coeficientes constantes son:
    \[
    y_1(t)=e^{\alpha t}\cos(\beta t),\qquad y_2(t)=e^{\alpha t}\sin(\beta t).
    \]
    Por lo tanto, la soluci\'on general es:
    \[
    y(t)=c_1e^{\alpha t}\cos(\beta t)+c_2e^{\alpha t}\sin(\beta t),
    \]
    donde $c_1,c_2$ son constantes reales.
\end{theorem}

\begin{proof}[Justificaci\'on]
    La funci\'on compleja $e^{(\alpha+i\beta)t}$ es soluci\'on de la ecuaci\'on, pues $\alpha+i\beta$ es ra\'iz de la ecuaci\'on auxiliar. Por la f\'ormula de Euler,
    \[
    e^{(\alpha+i\beta)t}=e^{\alpha t}\cos(\beta t)+i\,e^{\alpha t}\sin(\beta t).
    \]
    Como los coeficientes $a,b,c$ son reales, al reemplazar en la ecuaci\'on la parte real y la parte imaginaria deben anularse por separado; por lo tanto, $e^{\alpha t}\cos(\beta t)$ y $e^{\alpha t}\sin(\beta t)$ son soluciones reales. Su Wronskiano es $\beta e^{2\alpha t}\neq 0$, as\'i que son L.I.
\end{proof}

\begin{example}
    Determine la soluci\'on del problema de valor inicial
    \[
    y''+2y'+2y=0,\qquad y(0)=0,\qquad y'(0)=2.
    \]
    La respectiva ecuaci\'on auxiliar es $r^2+2r+2=0$, la cual tiene por soluciones
    \[
    r=-1\pm i\Longrightarrow \alpha=-1,\quad \beta=1.
    \]
    Con lo cual, la soluci\'on general es:
    \[
    y(t)=c_1e^{-t}\cos(t)+c_2e^{-t}\sin(t).
    \]
    Luego debemos calcular las respectivas constantes. De $y(0)=0$ obtenemos $c_1=0$. Derivando,
    \[
    y'(t)=e^{-t}\big[(c_2-c_1)\cos(t)-(c_1+c_2)\sin(t)\big],
    \]
    y de $y'(0)=c_2-c_1=2$ obtenemos $c_2=2$. Por lo tanto, la soluci\'on del PVI es
    \[
    y(t)=2e^{-t}\sin(t).
    \]
\end{example}

\subsection*{Resumen}

\begin{center}
\renewcommand{\arraystretch}{1.4}
\begin{tabular}{|c|c|}
\hline
\textbf{Ra\'ices de } $ar^2+br+c=0$ & \textbf{Soluci\'on general} \\ \hline
$r_1\neq r_2$ reales & $c_1e^{r_1t}+c_2e^{r_2t}$ \\ \hline
$r_1=r_2=r$ real & $c_1e^{rt}+c_2te^{rt}$ \\ \hline
$\alpha\pm i\beta$ complejas & $e^{\alpha t}\big(c_1\cos(\beta t)+c_2\sin(\beta t)\big)$ \\ \hline
\end{tabular}
\end{center}

Estos casos se pueden extender f\'acilmente a ecuaciones de orden mayor que dos. Si bien no es el objetivo de este cap\'itulo, mostraremos un ejemplo para fijar la idea general.

\begin{example}
    Consideremos la ecuaci\'on:
    \[
    y'''+3y''-y'-3y=0,
    \]
    la cual tiene asociada la ecuaci\'on auxiliar $r^3+3r^2-r-3=0$. Agrupando t\'erminos, esta ecuaci\'on se puede factorizar como
    \[
    r^3+3r^2-r-3=r^2(r+3)-(r+3)=(r-1)(r+1)(r+3)=0,
    \]
    con lo cual la respectiva soluci\'on general es:
    \[
    y_h=c_1e^{t}+c_2e^{-t}+c_3e^{-3t}.
    \]
\end{example}

\newpage
\subsection*{Aplicaci\'on: el sistema masa-resorte}

\begin{remark}
    Como mencionamos anteriormente, la mec\'anica del modelo masa-resorte se rige por la segunda ley de Newton: la posici\'on de un objeto de masa $m$ en el tiempo $t$ est\'a dada por $y(t)$, que satisface la ecuaci\'on
    \[
    my''(t)+by'(t)+ky(t)=0,
    \]
    cuyos t\'erminos se pueden identificar f\'isicamente como
    \[
    \text{[inercia]}\,y''+\text{[amortiguaci\'on]}\,y'+\text{[rigidez]}\,y=0.
    \]
    En un sistema f\'isico usual se tiene $m>0$, $b\geq 0$ y $k>0$.
\end{remark}

\begin{example}
    Determine la ecuaci\'on de movimiento para el sistema masa-resorte con
    $m=36\,\mathrm{kg}$, $b=12\,\mathrm{kg/s}$ y $k=37\,\mathrm{kg/s^2}$, donde las condiciones iniciales son $y(0)=0{,}7\,\mathrm{m}$ e $y'(0)=0{,}1\,\mathrm{m/s}$.

    La ecuaci\'on es $36y''+12y'+37y=0$, con ecuaci\'on auxiliar $36r^2+12r+37=0$. Su discriminante es $144-4\cdot 36\cdot 37=-5184=-72^2$, luego
    \[
    r=\frac{-12\pm 72i}{72}=-\frac{1}{6}\pm i\Longrightarrow \alpha=-\frac{1}{6},\quad \beta=1,
    \]
    y la soluci\'on general es
    \[
    y(t)=c_1e^{-t/6}\cos(t)+c_2e^{-t/6}\sin(t).
    \]
    De $y(0)=0{,}7$ obtenemos $c_1=0{,}7$. Adem\'as,
    \[
    y'(0)=-\frac{c_1}{6}+c_2=0{,}1\Longrightarrow c_2=0{,}1+\frac{0{,}7}{6}=\frac{13}{60}.
    \]
    Por lo tanto, la ecuaci\'on de movimiento es
    \[
    y(t)=0{,}7\,e^{-t/6}\cos(t)+\frac{13}{60}\,e^{-t/6}\sin(t).
    \]
    El factor $e^{-t/6}$ hace que la amplitud de las oscilaciones decaiga con el tiempo: es un movimiento \textbf{subamortiguado}.

    \begin{center}
        \includegraphics[width=0.8\textwidth]{Capitulo_2/grafico_resorte_v2.png}
    \end{center}
\end{example}

\begin{example}
    Considere ahora la ecuaci\'on
    \[
    36y''-12y'+37y=0.
    \]
    Su ecuaci\'on auxiliar tiene ra\'ices $r=\frac{1}{6}\pm i$, por lo que su soluci\'on general es:
    \[
    y(t)=c_1e^{t/6}\cos(t)+c_2e^{t/6}\sin(t).
    \]
    Observe que el \'unico cambio respecto del ejemplo anterior es el signo del coeficiente de $y'$. Un coeficiente de amortiguaci\'on negativo no corresponde a un resorte usual (ser\'ia un sistema que \emph{inyecta} energ\'ia), y su efecto es que ahora el factor $e^{t/6}$ hace crecer la amplitud de las oscilaciones sin l\'imite, como se aprecia en el gr\'afico:
    \begin{center}
        \includegraphics[width=0.6\textwidth]{Capitulo_2/masa_resorte_2.png}
    \end{center}
\end{example}

\begin{example}
    Considere la ecuaci\'on
    \[
    y''+5y'-6y=0,
    \]
    que tiene por soluci\'on general:
    \[
    y_h=c_1e^{t}+c_2e^{-6t}.
    \]
    Si imponemos la condici\'on inicial $y(0)=1$ con una velocidad inicial $y'(0)=v_0$, obtenemos el sistema $c_1+c_2=1$, $c_1-6c_2=v_0$, cuya soluci\'on da
    \[
    y(t)=\dfrac{6+v_0}{7}e^{t}+\dfrac{1-v_0}{7}e^{-6t}.
    \]
    Observe que el t\'ermino $e^{t}$ crece sin l\'imite, mientras que $e^{-6t}$ tiende a cero. Por lo tanto, la soluci\'on tiende a cero cuando $t\to\infty$ \emph{\'unicamente} si $v_0=-6$; para cualquier otra velocidad inicial, la soluci\'on diverge a $+\infty$ (si $v_0>-6$) o a $-\infty$ (si $v_0<-6$). Esto se puede visualizar como:
    \begin{center}
        \includegraphics[width=0.7\textwidth]{Capitulo_2/masa_resorte_3.png}
    \end{center}
\end{example}
